% ============================================================
%  Capítulo 2 — Estimación Puntual
% ============================================================
\section{Estimación Puntual}

\subsection{Estimadores y sus propiedades}

\begin{defi}[Estimador]
  Un \textbf{estimador} de $\theta$ es un estadístico $\hat{\theta} =
  \hat{\theta}(X_1,\ldots,X_n)$ cuyo valor se toma como aproximación del
  parámetro desconocido $\theta$.
\end{defi}

\subsubsection{Insesgamiento}

\begin{defi}[Estimador insesgado]
  $\hat{\theta}$ es \textbf{insesgado} (o no sesgado) para $\theta$ si
  \[
    \E_\theta[\hat{\theta}] = \theta \quad \forall\, \theta \in \Theta.
  \]
  El \textbf{sesgo} se define como $\Bias(\hat{\theta}) = \E[\hat{\theta}] - \theta$.
\end{defi}

\begin{ejemplo}
  Si $X_1,\ldots,X_n \iid (\mu, \sigma^2)$, entonces $\bar{X}$ es insesgado
  para $\mu$ y $S^2 = \tfrac{1}{n-1}\sum_{i=1}^n (X_i - \bar{X})^2$ es
  insesgado para $\sigma^2$.
\end{ejemplo}

\subsubsection{Error cuadrático medio}

\begin{defi}[ECM]
  El \textbf{error cuadrático medio} de $\hat{\theta}$ es
  \[
    \mathrm{ECM}(\hat{\theta}) = \E\bigl[(\hat{\theta} - \theta)^2\bigr]
    = \Var(\hat{\theta}) + [\Bias(\hat{\theta})]^2.
  \]
\end{defi}

\subsubsection{Consistencia}

\begin{defi}[Consistencia]
  $\hat{\theta}_n$ es \textbf{consistente} para $\theta$ si
  $\hat{\theta}_n \plim \theta$, es decir,
  \[
    \forall\,\varepsilon > 0: \quad
    P_\theta(|\hat{\theta}_n - \theta| > \varepsilon) \to 0
    \quad \text{cuando } n \to \infty.
  \]
\end{defi}

\begin{resultado}
  Si $\mathrm{ECM}(\hat{\theta}_n) \to 0$ cuando $n \to \infty$, entonces
  $\hat{\theta}_n$ es consistente para $\theta$.
\end{resultado}

\subsection{Métodos de estimación}

\subsubsection{Método de los momentos}

Igualar los $k$ primeros momentos poblacionales con los $k$ primeros momentos
muestrales y resolver el sistema de ecuaciones resultante en $\theta$.

\subsubsection{Máxima verosimilitud}

\begin{defi}[Función de verosimilitud]
  Dada una muestra observada $\mathbf{x} = (x_1,\ldots,x_n)$, la
  \textbf{función de verosimilitud} es
  \[
    L(\theta; \mathbf{x}) = \prod_{i=1}^n f(x_i; \theta), \quad \theta \in \Theta.
  \]
  El \textbf{estimador de máxima verosimilitud (EMV)} es
  $\hat{\theta}_{\text{MLE}} = \arg\max_{\theta} L(\theta;\mathbf{x})$.
\end{defi}

\begin{obs}
  Normalmente se trabaja con la \emph{log-verosimilitud}
  $\ell(\theta) = \log L(\theta)$ por conveniencia numérica.
\end{obs}

\begin{teo}[Propiedades asintóticas del EMV]
  Bajo condiciones de regularidad, si $n \to \infty$:
  \[
    \sqrt{n}(\hat{\theta}_{\text{MLE}} - \theta)
    \dlim \N\!\bigl(0,\, [I(\theta)]^{-1}\bigr),
  \]
  donde $I(\theta)$ es la \textbf{información de Fisher} de una observación.
\end{teo}
